% English translation of checked_pages/esq_119_022.tex; source: cote 119, site 22, PDF page 23, typescript page 20.
% AI-translated by GPT-6 Astra (OpenAI, ChatGPT Pro mode).
\scanpage{119/22}
\mspage{20}

of a combinatorial map retains a meaning over a base field, and
even over an arbitrary base ring. It also retains a
meaning for regular polyhedra of any dimension, by
replacing the cartographic group $\underline{C}_2$ with a suitable n-dimensional
variant $\underline{C}_n$. The case n = 1, i.e. the theory of
regular polygons in any characteristic, was the subject
of a DEA course in 1977/78, and already brought out some
new phenomena, as well as the usefulness of working in a \uline{projective}
space rather than an affine ambient space (here the affine plane).
This is due in particular to the fact that in certain
characteristics (and especially in characteristic 2) the
center of a regular polyhedron is thrust out to infinity. On the other
hand, the projective context, unlike the affine context,
allows one to develop with ease a duality formalism for
regular polyhedra, corresponding to the duality formalism
of combinatorial or topological maps (where the roles of vertices
and faces, in the case n = 2, say, are interchanged).
It turns out that for every projective regular polyhedron, one can
define an associated canonical hyperplane, which plays the role of
a canonical hyperplane at infinity and allows one to regard the given
polyhedron as an affine regular polyhedron.

The extension of the theory of regular polyhedra (and,
more generally, of all sorts of geometric-combinatorial
configurations, including root systems...) from the base field
$\mathbb{R}$ or $\mathbb{C}$ to a general base ring seems to me to have,
in this part of geometry, a scope comparable to the analogous
extension that has taken place since the beginning of the century in algebraic
geometry, or over the last twenty years or so in topology ( *),
with the introduction of the language of schemes and that of topoi.

\rule{4cm}{0.4pt}

(*) In writing this, I am aware that few topologists,
even today, recognize this conceptual and technical
broadening of topology and the resources
it offers.

\margin{22}

